\begin{table}[htbp]
\centering
\caption{Monthly labor income by educational attainment}
\label{tab:latam_country_education_incomes}
\small
\resizebox{\textwidth}{!}{%
\begin{tabular}{lrrrrrr}
\toprule
Economy & Low initial & Low final & Middle initial & Middle final & High initial & High final \\
\midrule
Argentina & 863 & 622 & 1,214 & 834 & 1,995 & 1,410 \\
Bolivia & 801 & 760 & 1,159 & 1,031 & 1,588 & 1,549 \\
Brazil & 523 & 535 & 767 & 759 & 2,125 & 1,944 \\
Chile & 632 & 674 & 1,080 & 1,019 & 2,660 & 2,337 \\
Colombia & 459 & 440 & 688 & 659 & 1,897 & 1,743 \\
Costa Rica & 855 & 838 & 1,374 & 1,259 & 2,951 & 2,681 \\
Dominican Republic & 640 & 648 & 726 & 797 & 1,628 & 1,533 \\
Ecuador & 577 & 501 & 790 & 713 & 1,604 & 1,269 \\
El Salvador & 465 & 587 & 773 & 849 & 1,208 & 1,303 \\
Guatemala & 454 & 535 & 907 & 763 & 1,526 & 1,556 \\
Mexico & 492 & 544 & 652 & 668 & 1,174 & 1,122 \\
Paraguay & 751 & 787 & 1,062 & 995 & 1,950 & 1,554 \\
Peru (Lima and Callao) & 548 & 485 & 699 & 644 & 1,507 & 1,240 \\
Uruguay & 861 & 786 & 1,237 & 1,204 & 2,216 & 2,254 \\
\bottomrule
\end{tabular}
}
\begin{minipage}{0.98\textwidth}
\footnotesize
\textit{Note:} Mean monthly labor income in 2017 PPP dollars. All observations correspond to the fourth quarter.
\end{minipage}
\end{table}
